\documentclass[Total.tex]{subfiles}
\begin{document}
		\chapter{Fundamental Groups}
	\section{Definitions}
	\begin{Def}
		Let $X$ be a topological space:
		\begin{itemize}
			\item We define a \textbf{path} in $X$ as a continuous map
			\begin{gather*}
				\gamma:[a,b]\rightarrow X
			\end{gather*}
			often we say: $\gamma(a)$ is the \textbf{starting point} and $\gamma(b)$ si the \textbf{end point}.
			\item A \textbf{loop} is a path with $\gamma(a)=\gamma(b)$;
			\item Given a path $\gamma:[a,b]\rightarrow X$, define its inverse 
			\begin{gather*}
				\bar{\gamma}:[a,b]\rightarrow X\qquad t\mapsto \gamma(b+a-t)
			\end{gather*}
		\end{itemize}
	\end{Def}
	
	\subsection{Path-Homotopy}
	\begin{Def}
		Let $X$ be a topological space:
		\begin{itemize}
			\item Let $\gamma_0,\gamma_1:[a,b]\rightarrow X$ be two paths with the same starting point $P$, and the same end point $Q$.
			
			A \textbf{path homotopy} between the paths $\gamma_0,\gamma_1$ is a homotopy rel $\{a,b\}$ between $\gamma_0$ and $\gamma_1$.
			
			In other words, a path-homotopy is a continuous map
			\begin{gather*}
				\Gamma:[a,b]\times [0,1]\rightarrow X
			\end{gather*}
			Such that
			\begin{gather*}
				\Gamma(t,0)=\gamma_0(t)\qquad \Gamma(t,1)=\gamma_1(t)\\
				\Gamma(a,s)=P\qquad \Gamma(b,s)=Q
			\end{gather*}
			\item If $\gamma_0$ and $\gamma_1$ are path-homotopic, then written as $\gamma_0\simeq_p \gamma_1$;
			\item Especially, a loop $\gamma$ is \textbf{null-homotopic}, if it is path-homotopic to the constant path $\delta$ given by
			\begin{gather*}
				\delta(t):=\gamma(a)
			\end{gather*}
		\end{itemize}
	\end{Def}
	\subsubsection{Reparametrization}
	\begin{Def}
		Let $\gamma:[a,b]\rightarrow X$ be a path in topological space $X$, and let
		\begin{gather*}
			\varphi:[a,b]\rightarrow [a,b]
		\end{gather*}
		which satisfies:
		\begin{itemize}
			\item Continuous;
			\item $\gamma(a)=a;\gamma(b)=b$;
			\item Monotone;
		\end{itemize}
		
		Then, $\gamma\circ\varphi$, which is called \textbf{reparametrization} of $\gamma$.
	\end{Def}
	\begin{lemma}
		(Reparametrization-Homotopy Lemma) Let $\gamma:[a,b]\rightarrow X$ be a path in a topological space $X$.
		
		If: a path $\delta:[a,b]\rightarrow X$ is obtained from $\gamma$ by a reparametrization, then
		\begin{gather*}
			\gamma\simeq_p \delta
		\end{gather*}
	\end{lemma}
	\begin{proof}
		Suppose $\varphi:[a,b]\rightarrow [a,b]$ a reparametrization such that
		\begin{gather*}
			\delta=\gamma\circ \varphi
		\end{gather*}
		
		Then, define the homotopy
		\begin{gather*}
			P:[a,b]\times [0,1]\rightarrow X\\
			(t,s)\mapsto \gamma(t+(1-s)\varphi(t))
		\end{gather*}
		
		is a path homomtopy
	\end{proof}
	\subsubsection{Path-Homotopy Class}
	
	\begin{lemma}
		(Path-Homotopy Equivalence Lemma) Let $X$ be topological space and $x,y\in X$. Then,path-homotopy is an equivalence relation on the set of paths $[a,b]\rightarrow X$.
	\end{lemma}
	\begin{proof}
		Path-homotopy satisfies the relation condition:
		\begin{itemize}
			\item $\gamma\sim_{p} \gamma$;
			\item Transitivity: Suppose $P:\gamma_1\sim_p \gamma_2,P':\gamma_2\sim_p \gamma_3$, then
			\begin{gather*}
				P''(s,t):=\begin{cases}
					P(s,2t)&t\in [0,\frac{1}{2}]\\
					P''(s,2t-1)&t\in [\frac{1}{2},1]
				\end{cases}
			\end{gather*}
			\item Symmetry: Omitted.
		\end{itemize}
	\end{proof}
	
	\begin{proposition}
		内容...
	\end{proposition}
	\begin{proof}
		内容...
	\end{proof}
	
	\begin{Def}
		\begin{itemize}
			\item (Path-Homotopy Class) Suppose $f:[a,b]\rightarrow X$ be path, denote $[f]$ the \textbf{path-homotopy class};
			\item (Product) Define
			\begin{gather*}
				\alpha*\beta(t):=\begin{cases}
					\alpha(2t)&t\in [0,\frac{1}{2}]\\
					\beta(2t-1)&t\in [\frac{1}{2},1]
				\end{cases}
			\end{gather*}
			
			\textbf{Remark} The product does not satisfy the associativity.
		\end{itemize}
	\end{Def}
	\begin{proposition}
		(Path-Homotopy Product Proposition) Suppose $X$ be a topological space, with paths $\alpha,\beta,\gamma:[0,1]\rightarrow X$. Then:
		
		\begin{itemize}
			\item If $\alpha(1)=\beta(0)$, $\beta(1)=\gamma(0)$, then
			\begin{gather*}
				[\alpha]\cdot ([\beta]\cdot [\gamma])=([\alpha]\cdot [\beta])\cdot [\gamma]
			\end{gather*}
			\item For $x\in X$, and $e_x:t\mapsto x$ is the constant path, the following equality holds:
			\begin{gather*}
				[e_{\alpha(0)}]\cdot [\alpha]=[\alpha]\cdot [e_{\alpha(1)}]=[\alpha]
			\end{gather*}
			\item The inverse path $\bar{\alpha}:t\mapsto \alpha(1-t)$, then
			\begin{gather*}
				[\alpha]\cdot [\bar{\alpha}]=[e_{\alpha(0)}]\qquad [\bar{\alpha}]\cdot [\alpha]=[e_{\alpha(1)}]
			\end{gather*}
		\end{itemize}
	\end{proposition}
	\begin{proof}
		\begin{itemize}
			\item The homotopy $\alpha*(\beta*\gamma)\simeq_p (\alpha*\beta)*\gamma$ is given by reparametrization
			\begin{gather*}
				\varphi:t\mapsto \begin{cases}
					\frac{1}{2}t&t\in [0,\frac{1}{2}]\\
					t-\frac{1}{4}&t\in [\frac{1}{2},\frac{3}{4}]\\
					2(t-\frac{3}{4})+\frac{1}{2}&t\in [\frac{3}{4},1]
				\end{cases}
			\end{gather*}
			\item Omitted;
			\item Omitted;
		\end{itemize}
	\end{proof}
	\subsubsection{The Fundamental Group of Pointed Topological Space}
	\begin{Def}
		Let $X$ be a topological space, and $x_0\in X$ a point.
		
		A \textbf{loop} is a path $f:[0,1]\rightarrow X$ such that
		\begin{gather*}
			f(0)=f(1)=x_0
		\end{gather*}
	\end{Def}
	\begin{proposition}
		($\pi_1$-Group Proposition) Let $X$ be a topological space and $x_0\in X$, the set
		\begin{gather*}
			\pi_1(X,x_0):=\{\mbox{path-homotopy classes of loops in }(X,x_0)\}
		\end{gather*}
		together with product $[\alpha*\beta]=[\alpha]\cdot [\beta]$ and unit $[e_{x_0}]$, which is a group.
	\end{proposition}
	\begin{proof}
		Define
		\begin{gather*}
			\pi_1(X,x_0)\times \pi_1(X,x_0)\rightarrow \pi_1(X,x_0)\\
			([\alpha],[\beta])\mapsto [\alpha]\cdot [\beta]=[\alpha*\beta]
		\end{gather*}
		
		Apply the homotopy-product proposition above.
	\end{proof}
	
	\begin{Def}
		Define $\pi_1(X,x_0)$ the \textbf{fundamental group} $X$ with base point $x_0$.
	\end{Def}
	\subsubsection{Change of Base Points, Type of Fundamental Groups}	
	\begin{proposition}
		(Change of Base Point Propsition) Let $X$ be a topological space, and $x_0,x_1\in X$, with $p:[0,1]\rightarrow X$ such that $p(0)=x_0,p(1)=x_1$.
		
		\begin{itemize}
			\item Map
			\begin{gather*}
				p_*:\pi_1(X,x_1)\rightarrow \pi_1(X,x_0)
			\end{gather*}
			
			is well-defined;
			\item If $x_0=x_1$, then $p_*$ is the conjugation by 
			\begin{gather*}
				p_*(g)=[p]\cdot g\cdot [p^{-1}]\in \pi_1(X,x_0),g\in \pi_1(X,x_1)
			\end{gather*}
			\item For $q:[0,1]\rightarrow X$ from $x_1$ to $x_2$ we have
			\begin{gather*}
				(p*g)_*=p_*\circ q_*:\pi_1(X,x_0)\rightarrow \pi_1(X,x_2)
			\end{gather*}
			\item If: $\pi_1(X,x_0)$ is \textbf{abelian}, then for any two paths $p,q:[0,1]\rightarrow X$ we have
			\begin{gather*}
				p_*=q_*:\pi_1(X,x_1)\rightarrow \pi_1(X,x_0)
			\end{gather*}
			
			And then, $p_*,q_*$ can be seen as action of conjugation, 
			\begin{gather*}
				p_*(g)=g=q_*(g)
			\end{gather*}
			
			Thus, $p_*=q_*$.
		\end{itemize}
	\end{proposition}
	\begin{proof}
		\begin{itemize}
			\item (1) (2) (3) Omitted;
			\item If $\pi_1(X,x_0)$ is abelian, then 
			\begin{gather*}
				p_*,q_*\in Iso(\pi_1(X,x_0),\pi_1(X,x_1))\cong Aut(\pi_1(X,x_0))
			\end{gather*}
		\end{itemize}
	\end{proof}
	
	\begin{corollary}
		(Base Points Independence) Let $x_1,x_2\in X$, $X$ is path-connected, then $\pi_1(X,x_1)\cong \pi_2(X,x_2)$.
	\end{corollary}
	\begin{proof}
		Omitted.
	\end{proof}
	
	\begin{Def}
		Let $X$ be a topological space, then 
		\begin{itemize}
			\item $\pi_1(X)$ is the \textbf{isomorphism type} of the fundamental groups;
			\item If $\pi_1(X)=0$, then $X$ is \textbf{simply connected}.
		\end{itemize}
	\end{Def}
	\subsubsection{Remark 1: Commutativity of Fundamental Group}
		\begin{example}
			Consider $X=\RR^2-(D_1\cup D_2)$, where $D_1\cap D_2=\emptyset$. Then, 
			\begin{gather*}
				X\simeq_p S_1\vee S_1
			\end{gather*}
			
			Its fundamental group(by Van Kampen) $\pi_1(X,x_0)\cong F(a,b)$. This is nonabelian.
		\end{example}
		
		As a matter of fact, we have the \textbf{Hurewicz map}
		\begin{gather*}
			h_n:\pi_1(X,x_0)\rightarrow H_n(X,\ZZ)
		\end{gather*}
		
		Especially, when $n=1$: $h_1:\pi_1(X,x_0)\rightarrow H_1(X,\ZZ)$,
		\begin{gather*}
			h_1:\pi_1(X,x_0)\rightarrow \pi_1(X,x_0)^{ab}\rightarrow H_1(X,\mathbb{Z})
		\end{gather*}
		
		Because: $h_1([\alpha*\beta])=h_1([\alpha])+h_1([\beta])\Rightarrow h_1(aba^{-1}b^{-1})=0$
		
		\begin{example}
			Abelian: $X=\RR^2$.
		\end{example}
	\subsubsection{Examples}
	\begin{itemize}
		\item Fundamental Group of the circle:
		\begin{theorem}
			(Fundamental Group of Circle) We have $\pi_1(S^1,1)\cong \ZZ$.
		\end{theorem}
		\begin{proof}
			Omitted.
		\end{proof}
		\begin{lemma}
			(Lesbegue's Lemma) Let
			\begin{gather*}
				f:[0,1]^m\rightarrow X
			\end{gather*}
			
			be a continuous map. Let $\{V_i\}$ be an open cover of $X$. 
			
			Then, for every $a_1,\dots,a_m\in \{0,1,\dots,N-1\}$there exists an $N\in \NN$ such that
			\begin{gather*}
				f([\frac{a_1}{N},\frac{a_1+1}{N}]\times\dots\times [\frac{a_m}{N},\frac{a_m+1}{N}])\subseteq V_i
			\end{gather*}
		\end{lemma}
		\begin{proof}
			Here: 
		\end{proof}
		\begin{Def}
			\begin{itemize}
				\item Consider map $\Xi:\RR\rightarrow S^1$;
				\item Let $X$ be a topological space with $f:X\rightarrow S^1$ continuous.
				
				A \textbf{lift} of $f$ is a continuous map
				\begin{gather*}
					\tilde{f}:W\rightarrow \RR
				\end{gather*}
				such that $\Xi\circ \tilde{f}=f$. In other words, such that for all $w\in W$ we have
				\begin{gather*}
					f(w)=exp(i\cdot \tilde{f}(w))
				\end{gather*}
				\begin{center}
					\begin{tikzcd}
						&\RR\arrow[d,"p"]\\
						W\arrow[ru,"\tilde{f}"]\arrow[r,"f"]&S^1
					\end{tikzcd}
				\end{center}
			\end{itemize}
		\end{Def}
		
		\begin{example}
			(Lift) Consider the map $f:[0,1]\rightarrow S^1,t\mapsto exp(2\pi int)$. Then, consider the lift given by
			\begin{gather*}
				\tilde{f}:t\mapsto 2\pi in\cdot t
			\end{gather*}
			The lift is not unique, take $k=2\pi\cdot n$.
		\end{example}
		
		\begin{lemma}
			(Lift Uniqueness Lemma) Let $W$ be a connected topological space. Let $f:W\rightarrow S^1$ be a continuous map and $\tilde{\varphi},\tilde{\psi}:W\rightarrow \RR$ be two lifts of $f$.
			
			If $\tilde{\varphi},\tilde{\psi}$ agrees on some point $w_0\in W$, then they agree everywhere. 
		\end{lemma}
		\begin{proof}
			Consider the following map
			\begin{gather*}
				\theta:W\rightarrow \RR\qquad w\mapsto \theta(w):=\tilde{\varphi}(w)-\tilde{\psi}(w)
			\end{gather*}
			Then, by hypothesis, we have
			\begin{gather*}
				\theta(w_0)=0\qquad exp(i\theta(w))=f(w)\cdot f(w)^{-1}=1
			\end{gather*}
			
			This implies, that $s\mapsto \theta(w)$ maps $W$ into $2\pi\cdot \ZZ=\Xi^{-1}(1)$. $W$ is connected, so they agrees everywhere.
		\end{proof}
		\begin{lemma}
			(Lift-Extension Lemma) Let $X$ be a topological space, we suppose that we have a decomposition that
			\begin{gather*}
				X=A\cup B\qquad A,B\mbox{ are closed}
			\end{gather*}
			Now, let $f:X\rightarrow S^1$ be continuous, $\tilde{f}:A\rightarrow \RR$ be a lift of $f|A$. We suppose that the following statements hold:
			\begin{itemize}
				\item There exists a $P\in S^1$ with
				\begin{gather*}
					f(B)\subseteq S^1-\{P\}
				\end{gather*}
				\item The intersection $A\cap B$ is \textbf{connected} and \textbf{nonempty}.
			\end{itemize}
			Then, there exists a lift $\tilde{f}:X\rightarrow \RR$ of $f$ agrees with the given lift on $A$.
		\end{lemma}
		\begin{figure}[H]
			\centering
			\includegraphics[width=0.8\textwidth]{AT1.png}
		\end{figure}
		\begin{proof}
			Firstly, pick some $x_0\in A\cap B$, and $s\in \RR$ such that $\Xi(s)=P$. Then,
			\begin{gather*}
				f(x_0)\in f(A\cap B)\subseteq f(B)\subseteq S^1-\{P\}\Rightarrow \tilde{f}(x_0)\notin\Xi^{-1}(P) 
			\end{gather*}
			So, there exists some $m\in \ZZ$ such that
			\begin{gather*}
				\tilde{f}(x_0)\in (s+2\pi m,s+2\pi (m+1))
			\end{gather*}
			Then, restriction of $\Xi$ to $(s+2\pi m,s+2\pi(m+1))\rightarrow S^1-\{P\}$ is a homeomorphism. Denote $\Psi:=[\Xi|(s+2\pi m,2+2\pi(m+1))]^{-1}$. 
			
			Now, define
			\begin{gather*}
				\tilde{f}(x):=\begin{cases}
					\tilde{f}(x)&x\in A\\
					\Psi(f(x))&x\in B\Rightarrow f(x)\in S^1-\{P\}
				\end{cases}
			\end{gather*}
			The definition agrees on $A\cap B$: For $x\in A\cap B$,
			\begin{gather*}
				\Psi(f(x))=\tilde{f}(x),x\in A
			\end{gather*}
			
			Thus, the definition of $\tilde{f}$ can be extended to $A\cup B$. 
		\end{proof}
		
		\begin{lemma}
			(Lift from $S^1$ to $\RR$) Given any continuous map 
			\begin{gather*}
				f:[0,1]\rightarrow S^1
			\end{gather*}
			and any $\tilde{y}_0\in \RR$ with
			\begin{gather*}
				f(0)=exp(i\cdot \tilde{y}_0)=\Xi(\tilde{y}_0)
			\end{gather*}
			There exists a \textbf{unique lift} $\tilde{f}:[0,1]\rightarrow \RR$ of $f$ with
			\begin{gather*}
				\tilde{f}(0)=\tilde{y}_0
			\end{gather*} 
		\end{lemma}
		\begin{proof}
			\begin{itemize}
				\item Uniqueness: By Lift Uniqueness lemma.
				\item Existence: Apply Lift-Extension lemma. Decompose the interval $[0,1]$ into smaller intervals, such that: on each subinterval, the image of $f$misses at least one point in $S^1$.
			\end{itemize}
		\end{proof}
		\begin{theorem}
			The map
			\begin{gather*}
				\Psi:\ZZ\rightarrow \pi_1(S^1,1)\qquad \\
				\Theta:\pi_1(S^1,z_0)\rightarrow \ZZ\qquad 
			\end{gather*}
			are well defined, and are isomorphisms of group.
		\end{theorem}
		\begin{proof}
			
		\end{proof}
		\item Lifting of Square:
		\begin{lemma}
			内容...
		\end{lemma}
		\begin{proof}
			内容...
		\end{proof}
		\item Fundamental Group of spheres:
		\begin{lemma}
			(Sphere Minus Point Lemma)For every $P\in S^n$, $S^n-\{P\}$ is simply connected.
		\end{lemma}
		\begin{proof}
			We have $S^n-\{P\}\cong \RR^n$. 
		\end{proof}
		\begin{proposition}
			For $n\geq 2$ we have
			\begin{gather*}
				\pi_1(S^n)=0
			\end{gather*}
			i.e. it is simly connected.
		\end{proposition}
		\begin{proof}
			内容...
		\end{proof}
		
	\end{itemize}
	
	\begin{proposition}
		($\pi_1$-Group Proposition) Let $X$ be a topological space and $x_0\in X$, the set
		\begin{gather*}
			\pi_1(X,x_0):=
		\end{gather*}
		together with product
	\end{proposition}
	\begin{proof}
		内容...
	\end{proof}
	
	\begin{Def}
		内容...
	\end{Def}
	\subsubsection{Alternative Descriptions}
	
	\subsubsection{Change of Base Points}
	\subsection{Applications}
	\subsubsection{The Fundalmental Theorem of Algebra}
	\begin{theorem}
		(TFTA) Every nonconstant polynomial in $\CC[z]$ has a zero in $\CC$.
	\end{theorem}
	\begin{proof}
		内容...
	\end{proof}
	\subsubsection{2D-BOrsuk-Ulam Theorem}
	\subsubsection{Path Integrals, Fundamental Groups}
	\subsection{Functoriality}
	\subsubsection{Group Homomorphisms induced by Continuous Map}
	\begin{Def}
		内容...
	\end{Def}
	\section{Coverings and Fundamental Groups}
	\subsection{Group Action}
	\subsubsection{Klein Bottle as Quotient}
	\begin{lemma}
		内容...
	\end{lemma}
	\begin{proof}
		内容...
	\end{proof}
	\subsection{Covering of Topological Spaces}
	\begin{Def}
		Let $p:X\rightarrow B$ be a continuous map between topological spaces:
		\begin{itemize}
			\item An open subset $U\subseteq B$ is called \textbf{uniformly covered} if 
			\item We say that the map $p:X\rightarrow B$ is a covering, if $p$ is surjective and if for every $b\in B$, there exists an open neighborhood $U$ of $b$ which is uniformly covered;
			\item Sometimes $X$ is called a \textbf{covering space} $B$.
		\end{itemize}
	\end{Def}
	\begin{lemma}
		(Covering-Preimage Lemma) Let $p:X\rightarrow B$ be a covering of topological spaces. If $B$ is a connected topological space, then for any two point $P,Q\in B$ we have
		\begin{gather*}
			\#p^{-1}(P)=\#p^{-1}(Q)
		\end{gather*}
	\end{lemma}
	\begin{proof}
		内容...
	\end{proof}
	
	\begin{Def}
		\begin{itemize}
			\item 
			\item 
			\item 
		\end{itemize}
	\end{Def}
	\subsubsection{Action-Covering Proposition}
	\begin{proposition}
		内容...
	\end{proposition}
	\begin{proof}
		内容...
	\end{proof}
	\subsection{Basic Properties of Coverings}
	\begin{proposition}
		Let $p:X\rightarrow Y$ be a covering. 
		\begin{itemize}
			\item Let $\gamma$ be a path in $Y$ wiht $p(x)=\gamma(0)$. Then there exists a unique $\dot{\gamma}$ in $X$ such that
			\begin{gather*}
				p\circ \dots{\gamma}=\gamma
			\end{gather*}
			\item Let $H$ be the homotopy between te paths $\gamma_0,\gamma_1$, in $Y$ and $l_1:[0,1]\rightarrow X$, be such that
			\begin{gather*}
				p(h(1))=H(0,t)
			\end{gather*}
			then, there exists unique $\dot{H}:[0,1]\times [0,1]\rightarrow X$ such that
			\begin{gather*}
				\dot{H}(0,t)=l_1(t),p(\dot{h})
			\end{gather*}
		\end{itemize}
		
		\begin{corollary}
			Let $p:(X,x)\rightarrow (Y,y)$ covering of pointed spaces, then
			\begin{gather*}
				p_*:\pi_1(X,x)\rightarrow \pi_1(Y,y)
			\end{gather*}
			is injective.
		\end{corollary}
		
		\begin{corollary}
			Let $\gamma$ be a loop in $(Y,y)$ and $\dot{\gamma}$ a lift of $\gamma$ shifting at $x$, then $\dot{\gamma}(1)\in p^{-1}(x)$ depends only on the homotopy class of $\gamma$.
			
			This  defines a mapy
			\begin{gather*}
				\pi_1(Y,y)\rightarrow p^{-1}(y)
			\end{gather*}
			If $x$ in $p_*\pi(Y,y)\subseteq p()$ $x(1)=x$,
			\begin{gather*}
				\pi_1(X,y)/p_*\pi_1(X,x)\rightarrow p^{-1}(y)
			\end{gather*}
		\end{corollary}
		\begin{proof}
			内容...
		\end{proof}
		
		\begin{proposition}
			(Classification of Covering Space) Assume $(Y,y)$ connected, locally pat connected ,semilocally simplicial 
			One assignment
			\begin{gather*}
				(X,x)\stackrel{p}{\rightarrow} (Y,y)
			\end{gather*}
		\end{proposition}
		
		\begin{corollary}
			Any space $(Y,y)$ as above has a simply connected covering space, unique up to iso
		\end{corollary}
		
		\begin{Def}
			Let $p(:X,x)\rightarrow (Y,y)$ its deck transformation group $Aut(p)$ is the group of automorphisms of $p$;
		\end{Def}
		
		\begin{proposition}
			Let $N\subseteq Aut(p)$ be the character of $p_*\pi(X,x)$ then $Aut(p)\rightarrow p^{-1}(y),\varphi\mapsto \varphi(x)$,$\pi_1(Y,y)/p_*\pi_1(X,x)$
		\end{proposition}
		have the same image in $p^{-1}(y)$. And composition is an iso of groups.
	\end{proposition}
	\begin{Def}
		$p$ is normal if $N=\pi_1(x,y)$.
	\end{Def}
	
	\begin{corollary}
		Let $p$ be normal, then:
		\begin{itemize}
			\item $Aut(p)=\pi_1(Y,y)/\pi_1(X,x)$;
			\item The action of $Aut(p)$ on $p^{-1}(y)$ is transitive;
			\item $p$ induces a homeomorphism
			\begin{gather*}
				X/Aut(p)\cong Y
			\end{gather*}
		\end{itemize}
	\end{corollary}
	\begin{proof}
		内容...
	\end{proof}
	
	\begin{corollary}
		Let $p:(X,x)\rightarrow (Y,y)$ be universal covering, let $\gamma_0,\gamma_1$ loops in $(Y,y)$, TFAE:
		\begin{itemize}
			\item $X_0\sim X_1$;
			\item $\dot{\gamma_0}(1)=\dot{\gamma_1}(1)$;
		\end{itemize}
		Moreover, $\gamma\leftrightarrow \varphi$
	\end{corollary}
	\subsection{Lifts}
	\subsection{Pullbacks}
	\section{Group Actions, and Fundamental Groups}
	
	\chapter{The Fundamental Groupoids}
	\chapter{Homotopy}
		\section{Homotopy of Maps}
			\begin{Def}
				Suppose $X,Y$ be topological spaces.
				
				$f,g\in Map_{Top}(X,Y)$ are called  \textbf{homotopic} if there exists a continuous map
				\begin{gather*}
					H:I\rightarrow Map_{Top}(X,Y)\qquad H(0)=f\qquad H(1)=g
				\end{gather*}
				
				such $H$ is calle a \textbf{homotopy}. Write $f\simeq g$.
				
				\textbf{Remark} The homotopies of some unpointed space are sometimes called \textbf{free homotopies}.
				
				\textbf{Remark} Sometimes for unpointed space, we use a full subcategory $T_\circ$, with better properties.
			\end{Def}
			
			\begin{lemma}
				Homotopy is an equivalence relation.
			\end{lemma}
			\begin{proof}
				\begin{itemize}
					\item 
					\item 
					\item 
				\end{itemize}
			\end{proof}
			
			\begin{Def}
				Define:
				\begin{itemize}
					\item $<X,Y>:=\{$ Free homotopy classes $X\rightarrow Y\}$;
					\item $[X,Y]:=\{$ Pointed Classes of Maps $X\rightarrow Y\}$;
				\end{itemize}
			\end{Def}
			\subsection{Adjoint Definition}
				Equivalently:
				\begin{Def}
					\begin{itemize}
						\item $f,g:X\rightarrow Y$ for unpointed space,
						\begin{center}
							\begin{tikzcd}
								X\arrow[rd]\arrow[r,"\iota_0"]&X\times I\arrow[d]&X\arrow[ld]\arrow[l,"\iota_1"]\\
								&Y
							\end{tikzcd}
						\end{center}
						\item $f,g:X\rightarrow Y$ for pointed space,
						\begin{center}
							\begin{tikzcd}
								X\arrow[rd,"f"]\arrow[r]&X\rtimes I\arrow[d]&X\arrow[l]\arrow[ld]\\
								&Y
							\end{tikzcd}
						\end{center}
					\end{itemize}
				\end{Def}
				
				And
				\begin{proposition}
					\begin{itemize}
						\item Homotopy is equivalent to
						\item Free homotopy is equivalent to
					\end{itemize}
				\end{proposition}
				\begin{proof}
					内容...
				\end{proof}
			\subsection{Homotopies of Paths}
		\section{Constructing Homotopies}
			\subsection{Straight-Line Homotopy}
		\section{Homotopy Theory}
		
			\subsection{Homotopy Extension Property, HEP}
				\begin{Def}
					Consider a pair $(A,X)$ with continuous function $i:A\rightarrow X$, then it satisfies the \textbf{HEP} with respect to space $Y$ if:
					
					For any map $f:X\rightarrow Y$, and left homotopy $\eta:A\times [0,1]\rightarrow Y$ such that
					\begin{gather*}
						\eta(-,0)=f\circ i
					\end{gather*}
					
					There exists a left homomtopy $\hat{\eta}:X\times [0,1]\rightarrow Y$ such that
					\begin{gather*}
						\hat{\eta}\circ (i\times id)=\eta
					\end{gather*}
				\end{Def}
				
				In diagram:
				\begin{center}
					\begin{tikzcd}
						内容...
					\end{tikzcd}
				\end{center}
				
				such $i$ is called \textbf{Hurewicz cofibration}.
			\subsection{Whitehead Product}
				\begin{Def}
					It is actually a \textbf{graded commutator}.
				\end{Def}
		\section{Groups, Cogroups}
		\section{Homotopy Groups}
			\begin{Def}
				Define:
				\begin{gather*}
					\pi_n(X):=[S^n,X]
				\end{gather*}
			\end{Def}
			
			\begin{theorem}
				\begin{itemize}
					\item 
					
					\item 
				\end{itemize}
			\end{theorem}
			\begin{proof}
				内容...
			\end{proof}
			
			\subsection{Uniqueness of Lift}
		\section{Homotopy and Duality}
			\begin{Def}
				内容...
			\end{Def}
		\section{Relative Homotopy Groups}
		\section{Deformation, Contraction}
			\subsection{Deformation}	
			\begin{Def}
				Suppose $i:A\hookrightarrow X,r:X\rightarrow A$, such that
				\begin{gather*}
					r\circ i=id_A\qquad i\circ r\simeq id_X
				\end{gather*}
				
				Then, $X$ deformation retracts onto $A$. 
			\end{Def}
			
			Equivalnetly, $r|A=id_A$. 
			
			\begin{proposition}
				(Deformation Retraction by Homotopy) The followings are equivalent:
				\begin{itemize}
					\item There is a homotopy $F:[0,1]\times X\rightarrow X$ where:
					\begin{itemize}
						\item $F(0,-)=id_X$;
						\item  $F(t,-)|A=id_A$;
						\item $F_1(X)=A$;
					\end{itemize}
					\item $X$ has deformation retracts onto $A$
				\end{itemize}
			\end{proposition}
			\begin{proof}
				\begin{itemize}
					\item $\Rightarrow$ $r=F_1$. 
					\item $\Leftarrow$ Notice that, we obtain a homotopy
					\begin{gather*}
						G:[0,1]\times X\rightarrow X\qquad G_0=id_X\qquad G_1=i\circ r
					\end{gather*}
					
					Now, $G_0(a)=G_1(a)=a$, but this does not imples, that $G_t(a)=a$. Thus we will construct a homotopy of homotopies
					\begin{gather*}
						K:A\times I\times I\rightarrow X\qquad K(a,t,s):=G(a,(1-s)(1-|2t-1|))
					\end{gather*}
					
					from $G|A\times I$ to constant homotopy $(a,t)\mapsto a$. Finally, take 
					\begin{gather*}
						H:X\times I\times I\rightarrow X
					\end{gather*}
					
					be the extension, and $F(x,t)=H(x,t,1)$. We obtain that
					\begin{gather*}
						H_t|A=id_A
					\end{gather*}
				\end{itemize}
			\end{proof}
	\chapter{Applications}
		\section{Applied Topology}
		\section{Covering of Smooth and Topological Manifolds}
		\section{title}
		\section{Seifert van Kampen Theorem}
		Firstly, the notion
		\begin{gather*}
			A*B:=\mbox{Coproduct of A and B}\\
			A*_G B:=(A*B)/<<\{\alpha(g)*\beta(g)^{-1}\}_{g\in G}>>,\alpha:G\rightarrow A,\beta:G\rightarrow B
		\end{gather*}
		\begin{center}
			\begin{tikzcd}
				G\arrow[r]\arrow[d]&A\arrow[d]\arrow[rdd]\\
				B\arrow[r]\arrow[rrd]&A*_G B\arrow[rd,dash]\\
				&&W
			\end{tikzcd}
		\end{center}
		\begin{lemma}
			\begin{itemize}
				\item The group homomorphsm $\Xi:\pi_1(U,x_0)*\pi_1(V,x_0)$ is surjective for: Topological space $X$ and $U,V\subseteq X$ be open subsets such that
				\begin{itemize}
					\item $X=U\cup V$;
					\item $U\cap V$ is path-connected and non-empty;
				\end{itemize}
				We choose a base point $x_0\in U\cap V$, then $x_0\in U\cap V$. 
				\item Every loop $s:[0,1]\rightarrow X$ in $x_0$ is path-homotopic to to a loop which is the concatenation of finitely many loops, each of which lies in $(U,x_0)$ or in $(V,x_0)$.
			\end{itemize}
		\end{lemma}
		\begin{proof}
			Let $s:[0,1]\rightarrow X=U\cup V$ be a loop. Then, $U,V$ are open, apply Lesbegue numbmer lemma: There exists a subdivision
			\begin{gather*}
				0=t_0<t_1<\dots<t_{k+1}=1	
			\end{gather*}
			of the interval $[0,1]$ such that the image each $s_i=s|[t_i,t_{i+1}]$ lies in $U$ or in $V$. 
			
			Then, by taking union of these intervals, we can satisfy that
			\begin{gather*}
				s(t_i)\in U\cap V
			\end{gather*}
			But such $s_i$ may be not a loop. We could construct loop for each $s_i$ by following steps: Notice that $U\cap V$ is path-connected. So we could take another path $s(t_{i+1})\rightarrow s_i$ in $U\cap V$. So we have
			\begin{gather*}
				s\simeq_p s_0*\dots*s_k\simeq_p (s_0*p_1)*(\bar{p_1}*s_1*p_2)*\dots 
			\end{gather*}
		\end{proof}
		\begin{theorem}
			(Seifert-van-Kampen Theorem) Let $X$ be a topological space and $X=U\cup V$ be a decomoposition of $X$ in two open subsets $U$ and $V$ such that: $U\cap V$ is 
			\begin{itemize}
				\item nonempty;
				\item path-connected
			\end{itemize}
			Let $x_0\in U\cap V$. We denote $i:U\cap V\rightarrow U$ and $j:U\cap V\rightarrow V$ the inclusions:
			\begin{itemize}
				\item The inclusions $k:U\rightarrow X$ and $l:V\rightarrow X$ induce an isomorphism
				\begin{gather*}
					\Phi:\pi_1(U,x_0)*\pi_1(V,x_0)\stackrel{\cong}{\rightarrow} \pi_1(X,x_0)
				\end{gather*}
				which has the property that the following diagram commutes
				\begin{center}
					\begin{tikzcd}
						\pi_1(U\cap V,x_0)\arrow[r,"i_*"]\arrow[d,"j_*"]&\pi_1(U,x_0)\arrow[rdd,bend left=10,"k_*"]\arrow[d]\\
						\pi_1(V,x_0)\arrow[r]\arrow[rrd,"l_*",bend right=10]&\pi_1(U,x_0)*_{\pi_1(U\cap V,x_0)} \pi_1(V,x_0)\arrow[rd,"\Phi"]\arrow[rd,"\cong",swap]\\
						&&\pi_1(X,x_0)
					\end{tikzcd}
				\end{center}
				group $\pi(X,x_0)$ is a pushout.
				\item Especially:If $U\cap V$ is simply connected, then the inclusions $U\rightarrow V$ and $V\rightarrow X$ induce an isomorphism
				\begin{gather*}
					\Phi:\pi_1(U,x_0)*\pi_1(V,x_0)\rightarrow \pi_1(X,x_0)
				\end{gather*}
			\end{itemize}
		\end{theorem}
		\begin{proof}
			Notice that: There exists pushout $\Phi:\pi_1(U,x_0)*\pi_1(V,x_0)\rightarrow \pi_1(X,x_0)$:
			\begin{itemize}
				\item Surjective:Consider following diagram
				\begin{center}
					\begin{tikzcd}
						\pi_1(U,x_0)*\pi_1(V,x_0)\arrow[rd]\arrow[d,"\Xi"]\\
						\pi_1(U,x_0)*_{\pi_1(U\cap V,x_0)} \pi_1(V,x_0)\arrow[r]&\pi_1(X,x_0)
					\end{tikzcd}
				\end{center}
				$\Xi$ given by the products. So surjective.
				\item Injective: $\Phi:\pi_1(U,x_0)*_{\pi_1(U\cap V,x_0)} \pi_1(V,x_0)\rightarrow \pi_1(X,x_0)$ which is induced by the inclusion $i:U\rightarrow X$ and $j:V\rightarrow X$ is injective.
				
				We just need to show: Exists map  $\Psi:\pi_1(X,x_0)\rightarrow \pi_1(U,x_0)*_{\pi_1(U\cap V,x_0)} \pi_1(V,x_0)$ such that 
				\begin{gather*}
					\Psi\circ \Phi=id
				\end{gather*}
			\end{itemize}
		\end{proof}
		\subsection{Suspension}
		\begin{Def}
			
		\end{Def}
		\begin{lemma}
			(Suspension-$\pi_1$-Lemma)
		\end{lemma}
		\subsection{Fundamental Groups of Wedges}
		\begin{Def}
			(Good Point) 
		\end{Def}
		
		\begin{lemma}
			内容...
		\end{lemma}
		\begin{proof}
			内容...
		\end{proof}
		\section{Classification of 2-dimensional Manifold}
		\section{Decision Problem}
		\section{Directed Colimits}
		\section{HNN-Seifert-Van Theorem}
		\section{Curves and Surfaces}
		\section{Introduction to Higher Homotopy Theory}
	\chapter{Weak Homotopy Equivalence}
		\section{Definition}	
			\begin{Def}
			\end{Def}
			\subsection{Serre Cofibration}
			\begin{Def}
				内容...
			\end{Def}
			\subsection{title}
			\subsection{Pushout and Cofibrations}
			\subsection{Properties}
		\section{CW-Approxiamation}
	\chapter{Cofibration, Fibrations} 
		\section{Cofibrations}
			\subsection{HEP Property}
				\begin{Def}
					\begin{itemize}
						\item 
						Given $(i,f,H_A)$ such that
						\begin{gather*}
							i:A\rightarrow X\qquad f:X\rightarrow Y\qquad H_A:A\times I\rightarrow Y
						\end{gather*}
						
						where $H_A$ is a homotopy . An \textbf{extension} of homotopy $H_A$ to $X$ is a homotopy $H:X\times I\rightarrow Y$ with the following diagram commutes:
						\begin{center}
							\begin{tikzcd}
								A\arrow[r,"i"]\arrow[d,"\iota_0"]&X\arrow[d,"\iota_0"]\arrow[ddr,bend left=20,"f"]\\
								A\times I\arrow[r,"i\times id_I"]\arrow[rrd,bend right=20,"H_A"]&X\times I\arrow[rd,dashrightarrow]\\
								&&Y
							\end{tikzcd}
						\end{center}
						
						Equivalently,
						\begin{center}
							\begin{tikzcd}
								A\arrow[r,"H_A"]\arrow[d,"i"]&Y^I\arrow[d,"ev_0"]\\
								X\arrow[r,"f"]\arrow[ru,dashrightarrow]&Y
							\end{tikzcd}
						\end{center}
						\item Homotopy Extension Propery: We say $i:A\rightarrow X$ has \textbf{homotopy extension property} \underline{with respect to} $f$;
						\item Especially, if $i$ has HEP for all $f$, then $i$ is called a \textbf{cofibration}.
					\end{itemize}
				\end{Def}
				
				\textbf{Remark} Sometimes, it is called \textbf{Hurewicz cofibration}
				
				\begin{proposition}
					The two definitions of extensions are equivalent.
				\end{proposition}
				\begin{proof}
					Just check the definition.
				\end{proof}
				
				\begin{proposition}
					The composition of coribrations is cofibration.
				\end{proposition}
				\begin{proof}
					Cosider the cofibrations $A\rightarrow B,B\rightarrow C$, then for any $f:C\rightarrow Y$,
					
					\begin{center}
						\begin{tikzcd}
							A\arrow[d,"\iota_0"]\arrow[r,"i"]&B\arrow[d,"\iota_0"]\arrow[rrdd,bend left=20]\arrow[r,"j"]&C\arrow[d,"\iota_0"]\arrow[rdd,bend left=20]\\
							A\times I\arrow[r]\arrow[rrrd,bend right=20]&B\times I\arrow[r]\arrow[rrd,bend right=20]\arrow[rrd,dashrightarrow]&C\times I\arrow[rd,dashrightarrow]\\
							&&&Y
						\end{tikzcd}
					\end{center}
				\end{proof}
				
				\begin{proposition}
					Suppose: The \underline{solid arrow} of the diagram is \textbf{homotopy commutative}. And $i$ is a cofibration,
					\begin{center}
						\begin{tikzcd}
							A\arrow[r,"i"]\arrow[d,"f"]&B\arrow[d,"\gamma",swap,dashrightarrow]\arrow[d,bend left=20,"g"]\\
							C\arrow[r,"j"]&D
						\end{tikzcd}
					\end{center}
					
					then, $g$ can be replaced with a homotopic map $\gamma:B\rightarrow D$ making the square strictly commutative.
				\end{proposition}
				\begin{proof}
					We just consider the simplified form as following.
				\end{proof}
				
				\begin{corollary}
					Especially,
					\begin{center}
						\begin{tikzcd}
							A\arrow[r]\arrow[rd]&B\arrow[d,dashrightarrow]\arrow[d,bend left=20]\\
							&D
						\end{tikzcd}
					\end{center}
				\end{corollary}
				\begin{proof}
					Suppose we have
					\begin{gather*}
						H:A\times I\rightarrow D
					\end{gather*}
					Then, consider the HEP
					\begin{center}
						\begin{tikzcd}
							A\arrow[r,"i"]\arrow[d]&B\arrow[d]\arrow[rdd]\\
							A\times I\arrow[r]\arrow[rrd,"H"]&B\times I\arrow[rd,dashrightarrow]\\
							&&D
						\end{tikzcd}
					\end{center}
					
					Just take the value of new homotopy at $0,1$.
				\end{proof}
				
				\subsubsection{Pointed-Topology Problems}
				
				\begin{proposition}
					\begin{itemize}
						\item The inclusion $\emptyset\rightarrow X$ is a cofibration;
						\item Every homeomorphism is a cofibration;
						\item A cofibration $i:A\rightarrow X$ is embedding, i.e.
						\item If $i$ is an embedding, then $i$ is a cofibration$\Leftrightarrow i(A)\hookrightarrow X$ is a cofibration;
						\item If $i:A\hookrightarrow X$ is a cofibration, then there is an open set $U\subseteq X$ containing $A$ such that: $A$ is a retract of $U$;
						\item If $i:A\hookrightarrow X$ is a cofibration, then $A$ is a closed subset of $X$;
						\item Inclusion
						\begin{gather*}
							\{0\}\hookrightarrow \{0\}\cup \{\frac{1}{n}|n\in \ZZ_+\}\qquad \QQ\hookrightarrow \RR
						\end{gather*}
						
						are not cofibrations.
					\end{itemize}
				\end{proposition}
				\begin{proof}
					\begin{itemize}
						\item 
						\item Omitted;
						\item 
						\item 
						\item 
					\end{itemize}
				\end{proof}
				
				\subsection{Pushout}
					\begin{Def}
						A \textbf{pushout square} is a square
						\begin{center}
							\begin{tikzcd}
								内容...
							\end{tikzcd}
						\end{center}
					\end{Def}
					\begin{proposition}
						$i$ is a cofibration$\Leftrightarrow$ The identity $id_T$ can be extended to a homotopy $X\times I\rightarrow T$;
					\end{proposition}
					\begin{proof}
						内容...
					\end{proof}
					
					\begin{proposition}
						内容...
					\end{proposition}
					\begin{proof}
						内容...
					\end{proof}
					
					\begin{corollary}
						内容...
					\end{corollary}
					\begin{proof}
						内容...
					\end{proof}
				\subsubsection{$S^n,D^{n+1}$}
				\subsubsection{Embedding in }
				\subsubsection{Retraction}
				
				\begin{Def}
					内容...
				\end{Def}
			\subsection{NDR Square, Str$\emptyset$m Theorem}
				\begin{Def}
					(NDR Square)
				\end{Def}
			\subsection{Product}
			\subsection{Closed Cofibration}
			
				Especially
				\begin{Def}
					Suppose: $i:A\hookrightarrow X$ maps $A$ into $X$ as Homeomorphism, and $i(A)\subseteq X$ is closed. And $i$ has \textbf{HEP}.
				\end{Def}
		\section{Fibrations}
			\begin{Def}
				\begin{itemize}
					\item 
					We have the dual of object $Cyl\dashv P$, the dual of cofibration is
					\begin{center}
						\begin{tikzcd}
							X\arrow[rd,dashrightarrow,"K"]\arrow[rrd,bend left=20,"K_B"]\arrow[rdd,bend right=20,"f"]\\
							&PE\arrow[d,"ev_0"]\arrow[r]&PB\arrow[d,"ev_0"]\\
							&E\arrow[r,"p"]&B
						\end{tikzcd}
					\end{center}
					
					Given $(p,f,K_B)$ there exists a right homotopy $K$ making the diagram commute.
					
					\item (Homotopy Lifting Property,HLP)  $p$ has \textbf{homotopy lifting property} with respect to $f$.
					\item (Fibration) If holds for every $f$, then $p$ is a \textbf{fibration}.
				\end{itemize}
			\end{Def}
			
			\textbf{Remark} 'right homotopy' is used to describe the homotopy, which is used here. 
			
			Equivalently:
			\begin{center}
				\begin{tikzcd}
					X&E\\
					Cyl(X)&B
				\end{tikzcd}
			\end{center}
			
			It is actually a lift diagram.
			
			\begin{Def}
				\begin{itemize}
					\item 
					\item 
				\end{itemize}
			\end{Def}
			
			\begin{proposition}
				Suppose that the solid arrow part of square
				\begin{center}
					\begin{tikzcd}
						X&E\\
						Y&B
					\end{tikzcd}
					
					If $p:E\rightarrow B$ is a fibration, then $f$ can be replaced with a homotopic map $X\rightarrow E$ making the square strictly commutative.
				\end{center}
			\end{proposition}
			\begin{proof}
				内容...
			\end{proof}
			
			\subsection{Lifting Functions}
			\subsection{Pullbacks}
			\subsection{Serre Fibration}
			\begin{Def}
				\begin{itemize}
					\item Write
					\begin{gather*}
						J_{Top}:=\{D^n\stackrel{(id,\delta_0)}{\hookrightarrow} D^n\times I\}_{n\in \NN}\subseteq Mor(Top)
					\end{gather*}
					\item A \textbf{Serre fibration} is $p:X\rightarrow Y$ which is $J_{Top}$-injective. i.e. for any $(id,\delta_0)$, it has right-lifting property: For any diagram
					\begin{center}
						\begin{tikzcd}
							D^n&X\\
							D^n&I
						\end{tikzcd}
					\end{center}
				\end{itemize}
			\end{Def}
			
			
			\begin{proposition}
				内容...
			\end{proposition}
			\begin{proof}
				内容...
			\end{proof}
			\subsubsection{Hurewicz Fibration}
			\begin{proposition}
				Serre fibration is in particular a Hurewicz fibration;
			\end{proposition}
			\begin{proof}
				内容...
			\end{proof}
			\subsubsection{Kan Fibration}
			\begin{proposition}
				内容...
			\end{proposition}
			\begin{proof}
				内容...
			\end{proof}
		\section{Factoring through Cofibrations and Fibrations}
		\section{Fundmanetal Lifting Property} 
	\chapter{Obstruction Theory}
	\chapter{Stable Homotopy I}
		\section{Sequential Spectra, Stable Homotopy Group}	
			\subsection{Sequential Spectra}
			
			\subsection{Stable Homotopy Group}
\end{document}